\documentclass[11pt,letterpaper]{article}
\usepackage[T1]{fontenc}
\usepackage{lmodern,amsmath,amssymb,amsthm,microtype}
\usepackage[margin=0.9in]{geometry}
\usepackage[hidelinks]{hyperref}
\title{Adding any six points to G449 preserves five-colorability}
\author{Aafi Khan Malek\\ \small Independent researcher, Ahmedabad, India\\
\small\href{mailto:aafimalek2032@gmail.com}{aafimalek2032@gmail.com}}
\date{28 November 2026}
\begin{document}
\maketitle
\begin{abstract}
For the fixed 449-point unit-distance graph of Heule and
Przybocki, prior exact searches show five-colorability after the
addition of any five points, and after six points in several
cases. We enumerate possible components for the remaining grid
case and compare their compatible colorings of the base graph.
Exact finite replay verifies all ten disconnected component-size
partitions. Thus adding any six distinct points anywhere in the plane
to G449 preserves five-colorability. This does not settle any
seven-point extension or the chromatic number of the plane.
\end{abstract}

\section{Reduction to a finite component problem}
Let \(G\) be the unit-distance graph on the fixed
\href{https://github.com/bprzybocki/G449/tree/125bfda97ed1d336f84694de04dedcfb349c5164}{G449 points}
of \href{https://arxiv.org/abs/2611.28518}{Heule and Przybocki}.
Include every unit edge when adding points. The previous
\href{https://github.com/Aafimalek/novel-discoveries/tree/234e0ce194fdfa544d4b551ccea9898b2239dc69/hadwiger-nelson/verification/g449/prime-seven-catalogue/five-point-extensions}{five-point theorem}
implies that any non-five-colorable six-point extension would
have total degree at least five at every added point. A point
of lower degree could be removed, then restored using an unused
color. The exact
\href{https://github.com/Aafimalek/novel-discoveries/tree/6376de4361b98f605326dd6275fa5f9e8b5dd3fb/hadwiger-nelson/verification/g449/prime-seven-catalogue/six-off-grid-extensions}{off-grid}
and
\href{https://github.com/Aafimalek/novel-discoveries/tree/f7be4c235d3d6bc073058db7fb201663ed7955e7/hadwiger-nelson/verification/g449/prime-seven-catalogue/six-grid-one-nonuniversal}{low-original-degree}
results reduce a candidate to six points in
\(\tfrac17\mathbb Z[\zeta_{42}]\), each with at least two
neighbors in \(G\). The entire
\href{https://github.com/Aafimalek/novel-discoveries/tree/018eda729500e0bbafca0557172bd6b0afe0b337/hadwiger-nelson/verification/g449/prime-seven-catalogue/six-grid-connected}{connected case}
has already been covered, where connectivity refers to the
graph induced by the added points.

For the remaining disconnected case, the 546 exact signed
unit displacements give a graph on 7,056 candidate sites with
33,908 induced edges. Each candidate is found by translating
an original site by a displacement and keeping only outside sites
with at least two original neighbors. No floating point
predicate decides an edge. If \(C\) is a connected component
of the new-point induced graph, each \(v\in C\) must satisfy
\(b(v)+d_C(v)\ge5\), where \(b(v)\) is its number of original
neighbors and \(d_C(v)\) its degree inside the component; there
are no edges between different components.

The census enumerates all connected subsets of size at most
five using the canonical parent from the connected proof:
delete the largest nonroot vertex whose deletion leaves a
connected induced set, taking the least ID as root. Each set
has a unique path from its root. A partial set \(S\) can be
pruned only if some \(b(v)+d_S(v)+5-|S|<5\); no completion of
size at most five can then be eligible. The 64 root shards
give the size-one through size-five eligible counts
\[
 574,\qquad 2{,}079,\qquad 13{,}188,\qquad
 115{,}395,\qquad 1{,}119{,}839.
\]

\section{A common verified base coloring}
For each proper five-coloring of \(G\) in a finite bank,
the list of a new vertex consists of the colors absent on all
its original neighbors. For every enumerated component
\(C\), exact finite list coloring computes the acceptance mask
of all base colorings that extend over \(C\).
If the masks of the components have a common bit, the
extension is five-colorable: use the same base coloring and
combine the independent component extensions. Thus it suffices
to prove a nonempty bit intersection even for arbitrary masks
drawn from the respective component sizes. This relaxes the
geometry by ignoring collisions and cross-edges; a successful
universal check is still sufficient.

For a fixed size, it suffices to keep inclusion-minimal masks,
because every actual mask contains a minimal one.
We also form all intersections of two singleton masks
(unordered, with repetition), retaining only their inclusion-minimal
intersections. A bitset indexed by target masks records
the targets containing each palette bit. For a query mask,
the OR of the indexed bitsets for all its bits covers every
target precisely when no target mask is disjoint from the query.
This checks the necessary finite unions of products without
assuming geometry. Partial or capped enumeration makes no
whole-partition claim.

The 61-palette selection is entirely positive evidence:
all original color vectors are checked on the 2,198 original
edges and all resulting masks are regenerated. In particular,
one model on the union of the three earlier residual
configurations was independently checked on 2,325 unit edges of
466 distinct points, after recomputing 108,345 exact distances.
Its restriction to \(G\) and six exact rotations add seven
proper palettes. The exact 52-, 54-, and 61-palette profiles are
append-only, so earlier closures remain valid.

\section{Exact coverage and scope}
The 52-palette bank covers \((5,1)\), \((4,2)\), \((4,1,1)\), \((3,3)\), and \((2,2,2)\).
The 54-palette bank adds \((3,2,1)\) and \((2,2,1,1)\).
The 61-palette bank adds \((3,1,1,1)\) and \((2,1,1,1,1)\).
The last check uses \(610\) pair-component masks times
\(11{,}888\) minimal two-singleton intersections as queries:
all \(7{,}251{,}680\) have nonempty intersection with every
target two-singleton intersection. For six singletons, take two such
pair-intersection masks with repetition as a query, and use the
same family of pair-intersection masks as targets. There are
\(\binom{11{,}889}{2}=70{,}668{,}216\) unordered query pairs
with repetition. The exact join covers all of them and finds
no disjoint target. Every disconnected partition therefore
has a common palette. Together with the connected, off-grid,
and low-original-degree results, all configurations
of six new points are five-colorable. The earlier theorem covers
fewer than six. The exact enumeration and coverage replayers
are available with the source. This result is not checked in
Lean and does not improve the known 6--7 bound for the plane.

\medskip\noindent
This research note was developed with AI assistance.
\end{document}
